\section{Related Work}
\label{sec:related}

We organize prior work into three threads that converge on the gap this paper
addresses: the relationship between retrieval quality, context structure, and
generation faithfulness.

\subsection{RAG evaluation and faithfulness measurement}

The original RAG formulation coupled non-parametric retrieval with parametric
generation~\citep{lewis2020rag}. Subsequent architectures enhanced retrieval
through dense encoders~\citep{karpukhin2020dpr}, few-shot
learning~\citep{izacard2022atlas}, and hybrid strategies. On the evaluation
side, RAGAS~\citep{es2024ragas} and FActScore~\citep{min2023factscore}
established automated faithfulness assessment, while broader benchmarking
studies~\citep{chen2024benchmarking, gao2024retrievalaugmented} confirmed
that better retrieval is \emph{usually} correlated with better downstream
behavior.

The critical word is ``usually.'' These studies report average-case
correlations; they do not characterize the regime in which the
correlation breaks down. Our work targets precisely that regime — the
conditions under which retrieval improvements fail to transfer, or
actively harm, generation faithfulness.

\subsection{Retrieval-side interventions and context preparation}

Three lines of retrieval-side work are most relevant. First, chunking
strategies~\citep{kamradt2024chunking} determine the granularity of the
evidence pool but are typically tuned through retrieval accuracy rather than
faithfulness. Second, cross-encoder re-ranking~\citep{nogueira2019passage,
glass2022re2g} reorders candidates for relevance but does not consider how
the selected passages fit together. Third, context
compression~\citep{xu2023recomp} condenses evidence before generation,
reducing noise but operating on each passage independently.

\subsection{Adaptive and self-correcting RAG}

Recent work has moved quality control closer to the generator.
Self-RAG~\citep{asai2024selfrag} trains the model to emit reflection tokens
that decide whether to retrieve, whether retrieved passages are relevant, and
whether the generated output is supported---effectively learning to critique
its own retrieval. CRAG~\citep{yan2024crag} adds a corrective layer that
evaluates retrieval quality post-hoc: if the retrieved documents are judged
ambiguous or incorrect, the system triggers web search or knowledge
refinement before generation. RECOMP~\citep{xu2023recomp} compresses
retrieved passages into shorter summaries, reducing irrelevant content.

These methods represent important advances, but they share a common design
assumption: quality is assessed and optimized at the level of individual
passages or passage-generation pairs. Self-RAG's reflection tokens evaluate
whether a \emph{single} passage is relevant; CRAG's corrective action
decides whether the \emph{overall retrieval} is useful but does not examine
inter-passage consistency; RECOMP compresses each passage independently. None
of these methods directly measure or optimize the \emph{coherence among}
retrieved passages---whether the evidence set contains consistent
terminology, complementary details, and logical continuity across chunks.

Our work addresses an orthogonal failure mode. The refinement paradox we
document can occur even when every individual passage is relevant (as judged
by Self-RAG's reflection tokens or CRAG's quality evaluator), because the
problem lies not in passage-level relevance but in set-level coherence. A
retrieval set can consist entirely of individually relevant passages that
nonetheless fragment the context when assembled. HCPC-v2 and CCS are
therefore complementary to these methods: they could be integrated as an
additional coherence check within any of these pipelines, gating refinement
or re-retrieval on whether the assembled context is internally consistent
rather than only on whether individual passages are adequate.

\subsection{Context quality, distraction, and domain mismatch}

Several studies establish that context quality is more nuanced than relevance
alone. \citet{shi2023context} demonstrated that even partially irrelevant
context can degrade generation quality, showing that models are susceptible to
distraction from noisy passages. \citet{liu2024lost} revealed positional
bias in how language models use long contexts, with evidence in the middle
being systematically underweighted. Together, these findings suggest that the
\emph{structure} of the context---not just its relevance---matters for
generation.

On the domain side, \citet{thakur2021beir} showed that dense retrievers
trained on one distribution often fail on others, and
\citet{gu2024biomedical} demonstrated that domain-specific encoders are
required to capture specialized terminology. Our PubMedQA results connect
these retrieval-level observations to an answer-level consequence: when a
general-domain encoder produces low-coherence retrieval in a specialized
corpus, RAG can become less faithful than zero-shot prompting.

\subsection{Comparison with Adaptive and Self-Corrective RAG Methods}
\label{sec:baseline_compare}

Self-RAG~\citep{asai2024selfrag}, CRAG~\citep{yan2024crag}, and
RECOMP~\citep{xu2023recomp} are the most salient points of comparison.
All three optimize passage-level or passage-pair signals: Self-RAG's
reflection tokens gate retrieval and supportedness on individual
passages, CRAG's post-hoc classifier judges the overall retrieval as a
coarse binary, and RECOMP compresses each passage independently. None
explicitly scores the coherence of the assembled context as an object
read by the generator.

\paragraph{Why these methods would not fully resolve the paradox.}
The paradox occurs when every individual passage is locally relevant yet
the \emph{set} loses internal coherence through refinement, fragmentation,
or encoder-corpus mismatch. Self-RAG's tokens can still approve each
sub-chunk after splitting because each remains relevant in isolation; the
connective evidence lost between them is invisible at the passage level.
CRAG's binary quality verdict does not distinguish a coherent set from a
fragmented one. RECOMP's per-passage compression is the exact class of
operation the paradox predicts will harm faithfulness when applied
indiscriminately. All three are orthogonal to the coherence axis and can
co-occur with the paradox rather than prevent it.

\paragraph{Why we do not run a head-to-head empirical comparison.}
Our design isolates a single causal question---does preserving set-level
coherence recover faithfulness at matched retrieval similarity?---which
requires holding retrieval similarity and the generator fixed while
varying one structural property. Introducing Self-RAG or CRAG would
simultaneously change the retriever, the generator training, and the
decision policy, conflating the coherence effect with three other axes. A
system-comparison study answers a different question and is left to
follow-up work.

\paragraph{Integration, not replacement.}
A coherence gate (CCS-thresholded retrieval, or the HCPC-v2-style
selective refinement used here as a probe) is in principle insertable
inside Self-RAG's reflection step, before CRAG's corrective decision, or
between RECOMP's compression and generation. We conjecture such
integration reduces the paradox in each system and leave the measurement
to follow-up work.

\subsection{Lightweight Empirical Sanity Check with Adaptive RAG}
\label{sec:sanity_check}

We treat our rerank-only control (\S\ref{sec:method}) as a minimal
empirical sanity check against adaptive retrieval. It over-fetches
$k_{\text{fetch}}{=}7$ candidates, re-scores with a cross-encoder, and
returns the top~$k{=}3$---the same scoring-side mechanism CRAG relies on
before triggering corrective actions. Within our shared pipeline it
differs from HCPC-v1 in passage fragmentation and from HCPC-v2 in
coherence-aware gating, making it the cheapest ``retrieve-then-judge''
adaptive analogue compatible with our controlled design.

\paragraph{Expected behavior under the coherence account.}
The hypothesis makes three predictions before any numbers are consulted:
(i) rerank-only should not reproduce the refinement paradox, because it
reorders passages without fragmenting them; (ii) it should not close the
PubMedQA failure, because the coherence deficit there originates in the
encoder, not in ordering; (iii) it should track baseline faithfulness
more closely than HCPC-v1 and underperform HCPC-v2 in the regime where
selective intervention extracts coherence-preserving gains. Results in
Sections~\ref{sec:results} and~\ref{sec:analysis} are consistent with
each. This rules out the alternative explanation that the paradox is a
weakness specific to naive fixed-length retrieval: a scoring-based
``adaptive'' improvement, applied without coherence awareness, is
similarly insufficient.

\paragraph{Why we stop here.}
Extending to a full benchmark comparison would require training
Self-RAG's reflection-token model on our corpora, implementing CRAG's
web-augmented recovery loop, and matching both to our generator and
embedding stack---each adding confounds into a comparison meant to
isolate one structural property. The scoped sanity check strengthens
rather than replaces that future evaluation.

\subsection{Positioning of this work}

This paper sits at the intersection of RAG evaluation and retrieval
system design. We do not propose another retrieval score or reranking
architecture; Self-RAG, CRAG, and RECOMP each solve a different problem
(whether to retrieve, correcting bad retrievals, compressing passages),
and a head-to-head comparison would conflate multiple design axes and
obscure the variable we study: \emph{set-level coherence}. Instead, we
isolate coherence by varying only refinement policy within a shared
pipeline, enabling the controlled ablations in
Sections~\ref{sec:results} and~\ref{sec:analysis}. CCS and HCPC-v2 are
complementary to these methods and could be integrated as an additional
coherence gate within any of them.

Our contribution is the identification of a specific failure mode —
the refinement paradox — and its cause: a meso-level property of the
retrieved context, its coherence. This perspective unifies several
otherwise disconnected observations: why naive refinement fails on
domain-matched corpora, why selective refinement helps only under
specific conditions, why chunk size dominates prompt strategy, and why
domain-mismatched retrieval can make grounded generation less reliable
than ungrounded generation. \citet{ji2023hallucination} and
\citet{xu2024hallucination} argue that hallucination is an inherent
property of autoregressive generation; our results add that in RAG
settings, hallucination is also a property of the retrieval
configuration, with context coherence as the specific mediating
variable.
